\section{The component view, the code-level views, the sequence view, and the documentation approach}\label{sec:B}

\subsection{Component diagram (C4 Level 3): inside \texttt{fault-monitor}}\label{sec:B.1}

\begin{center}\includegraphics[height=0.72\textheight]{figures/component.pdf}\end{center}
\figcap{C4 Component diagram (Level 3), from \texttt{project/diagrams/component.d2}: inside
\texttt{fault-monitor}. Every box is one \emph{component}, tagged with that element type and
labelled with the exact type name so a reader can grep for it; where the source lives in
\texttt{shared/} rather than in \texttt{fault-monitor/} the box says so, because at runtime it is
a component of this container either way. The \texttt{boundary} group holds every framework and
transport type; \texttt{control} and \texttt{entity} depend on nothing outward.}

\subsection{Code diagrams (C4 Level 4), and why only two}\label{sec:B.2}

C4's own guidance is that the code level is \emph{optional} and worth drawing only for the parts
that warrant it; a class diagram of everything is a diagram nobody reads and that rots on the
first rename. The criterion used here is narrower than ``the complicated parts'': \textbf{a code
diagram earns its page where a claim this report makes is a claim about the shape of the code
rather than about the physics, and where the Level 3 box hides the very thing being claimed.} Two
places meet it. Everything else in the system is either a plain adapter, where the port name is
already the whole story, or logic whose interesting content is arithmetic that Sections \ref{sec:6} and \ref{sec:7}
give in full.

\begin{center}\includegraphics[width=\linewidth]{figures/code-diagnosis.pdf}\end{center}
\figcap{C4 Code diagram (Level 4), from \texttt{project/diagrams/code-diagnosis.d2}: the
diagnosis chain in \texttt{fault-monitor}'s \texttt{control} layer. Deliberately partial: only
the types an arrow lands on, and only the members that carry one of the claims below.
\texttt{MonitorParameters}' 25 settings (23 numbers and durations, plus the \texttt{lowerArm} and
\texttt{classifyWhen} mode enums), \texttt{OvernightFloor}, \texttt{AhuFlowModel} and
\texttt{FaultSignals}' twelve components stay at Level 3; \texttt{MonitorState}'s constants and
transitions are the state machine of Section \ref{sec:8.3} and are not drawn twice.}

Four claims in this report are claims about this figure, and at Level 3 all four are hidden
inside one box called \texttt{ZoneMonitor}.

\begin{itemize}
\tightlist
\item \textbf{The cycle has two phases and the order is load-bearing.}
  \texttt{FaultMonitor.runCycle()} calls \texttt{evaluate(now)} on all six monitors, then counts
  how many are \texttt{elevated()}, and only then calls \texttt{decide(signals, now)} on all six.
  The count is the common-mode signal that separates a blocked filter from a stuck damper, so a
  single-phase loop would hand room six a count mixed from this cycle and the last.
\item \textbf{The published residual is not the accumulated one.} \texttt{ZoneMonitor.evaluate}
  computes three quantities from one \texttt{Prediction}: \texttt{measured - predicted} goes out
  on \texttt{residuals/\{zone\}} because it is the honest number a human should see;
  \texttt{measured - ceiling} drives the upper CUSUM; and under the deployed \texttt{GUARDED}
  lower arm, \texttt{measured - floor} drives the lower one. The ceiling quantity is
  \emph{subtracted} rather than clipped at zero, which Section \ref{sec:7.3} explains is what keeps the
  in-control run length R-3 rests on.
\item \textbf{\texttt{serviced()} now has a transport, and the asymmetry this bullet used to
  describe is gone.} \texttt{FaultCase.serviced(ObservedAt)} is the only exit from
  \texttt{CONFIRMED} and \texttt{COMPENSATING}, and for most of this project's life nothing reached
  it: its only callers were two unit tests, so in the deployed system a false latch was permanent.
  It is reachable now by \texttt{ServiceReported} on \texttt{service/\{zone\}}, which
  \texttt{fault-monitor} subscribes to through a \texttt{ServiceRequests} adapter that validates it,
  and by \texttt{POST /api/zones/\{zone\}/cases/\{caseId\}/serviced} on the dashboard. The topic is
  deliberately \emph{not} under \texttt{commands/}: it carries a report about the world, not an
  instruction to an actuator, and putting it under \texttt{commands/} would have given a second
  publisher authority over a topic space the actuator owns. The full
  \texttt{CONFIRMED}~→~\texttt{COMPENSATING}~→~\texttt{SERVICED}~→~\texttt{WARMUP}~→~\texttt{HEALTHY}
  lifecycle is executable and tested end to end, including that a reported repair unwinds the
  compensation (Section \ref{sec:13.1}).
\item \textbf{\texttt{DAMPER\_STUCK\_\allowbreak OPEN} now exists, and naming it did not make it
  reachable.}
  \texttt{FaultClassi\allowbreak fication} holds \texttt{NONE}, \texttt{STUCK\_DAMPER},
  \texttt{DAMPER\_STUCK\_OPEN}, \texttt{BLOCKED\_FILTER}, \texttt{SENSOR\_DRIFT} and
  \texttt{UNCLASSIFIED}. Three things branch on the new constant: the classifier's low-side leaf,
  the supervisor's work-order narrative, and its severity. A damper stuck open is \texttt{WARNING}
  rather than \texttt{CRITICAL}, because that room is over-ventilated and safe while the other five
  are short. The two extremes are deliberately re-joined by \texttt{isUnresponsiveDamper()} for the
  one purpose where the direction is irrelevant: pinning the shared flow belief. \textbf{But the
  diagnosis is still unreachable at the deployed setting}, because only the unguarded lower CUSUM arm
  sees a damper stuck wide open and that arm costs six false confirmations per fault-free week
  against R-3's limit of one (Section \ref{sec:13}). So this makes the fault \emph{expressible and
  actionable}, not detectable, and the trade it sits behind is unchanged.
\end{itemize}

\begin{center}\includegraphics[height=0.72\textheight]{figures/code-compensation.pdf}\end{center}
\figcap{C4 Code diagram (Level 4), from
\texttt{project/diagrams/code-compensation.d2}: the compensation solver in
\texttt{supervisor}'s \texttt{control} layer. The inset lists the order of steps inside
\texttt{solve()}; none of them repeats. \texttt{PlantState},
\texttt{ProbeLimiter} and \texttt{WorkOrderNarrative} are Level 3, and the algebra itself is
Section \ref{sec:7.4} rather than this figure.}

The solver is the second case, and the reason is the same: Section \ref{sec:7.4}'s central claim is that
holding five healthy flows while multiplying the sixth is a \textbf{closed form rather than a
search}, and that is a statement about this code's shape that no box labelled ``the compensation
solver'' can carry. \texttt{CompensationSolver.solve} is straight-line: no loop over candidate
solutions, no convergence test, no per-room balance, and the only loop in the method writes the
result map. The three helper methods are each a closed form (\texttt{trimFloor},
\texttt{preservingSpeed}, \texttt{toleratedTrimFloor}), and the one conditional re-solve, when the
preserving speed would exceed the fan's ceiling, happens at most once. The two constants the
figure names are bounds rather than tuning knobs: \texttt{MAXIMUM\_REQUESTED\_GAIN} = 4.0 because
the dampers' 0.15 floor caps the trim at 0.25, and \texttt{HEALTHY\_FLOW\_TOLERANCE} = 0.05
because that is the most a healthy room may be asked to give up (or, read the other way, to be
over-supplied by, once there is no trim left to make).

\subsection{The key scenario in time (sequence diagram)}\label{sec:B.3}

The sequence view adds what the dynamic view leaves out: duration, and the three phases that
follow confirmation. Its fourth group is the point of the project: after compensation the
residual collapses to 20 ppm and the CUSUM to zero, and the case does \emph{not} clear.

\begin{center}\includegraphics[height=0.86\textheight]{figures/sequence.pdf}\end{center}
\figcap{Sequence diagram, from \texttt{project/diagrams/sequence.d2}. Four groups: the inner loop
filling the room; the outer loop confirming and responding; the latch holding after compensation
has made the CO₂ curve look healthy; and service, the only exit.}

\subsection{The documentation approach: C4 for the views, arc42 for the frame}\label{sec:B.4}

Two frameworks are used, for two different jobs, and each is weak where the other is strong.
\textbf{C4 supplies the structural views}: a numbered zoom from context to containers to components
to code, with a fixed element vocabulary at each level, so a reader can tell what kind of thing a box
is without reading its label twice. \textbf{arc42 supplies the document frame}: twelve numbered
sections that force the questions a structural diagram cannot answer (constraints, cross-cutting
concepts, quality scenarios, risks and technical debt). C4 has no place to put a measured quality
scenario or an honest debt register, and arc42 names a ``building block view'' without saying what a
building block may be.

The arc42 document is \texttt{project/docs/arc42.md}, and it deliberately does \emph{not} repeat
this report: where the report says something, arc42 cites the section and summarises it, and its
value is the two things this report has no natural home for: the quality tree with its measurable
scenarios (arc42 §10) and the debt register as of a stated date (arc42 §11). The C4 diagrams are
shared rather than duplicated: both documents reference the same nine files in
\texttt{project/diagrams/}.
\label{sec:Bend}
